\documentclass[11pt]{article}
\usepackage[margin=1in]{geometry}
\usepackage{amsmath,amssymb}
\usepackage{hyperref}
\usepackage{listings}
\usepackage{xcolor}
\usepackage{enumitem}

\lstset{
  basicstyle=\ttfamily\small,
  breaklines=true,
  frame=single,
  columns=fullflexible,
  keepspaces=true,
  showstringspaces=false
}

\title{AuditLane: A Reproducible Pre-Audit Cleanroom for the\\ Arbitrum Audit Program}
\author{Dipankar Sarkar \\ \small doom2quake}
\date{August 2026}

\begin{document}
\maketitle

\begin{abstract}
A project that lands audit funding still hands the auditor a raw repository: a build that
works on one machine and not another, a test suite that needs undocumented setup, and
deployment scripts that point at stale addresses. Under a subsidised audit program the
firm spends the first days of a fixed engagement reconstructing the build before a single
line is reviewed, and that reconstruction is paid for out of a subsidy meant to buy
review. We present AuditLane, a reproducible pre-audit cleanroom for the Arbitrum Audit
Program~\cite{arbaudit}. AuditLane folds a target Arbitrum repository into a single
content-addressed \emph{manifest}: a pinned toolchain, the captured test counts, and every
contract's source hash and Arbitrum Sepolia deployed-bytecode hash checked against the
locally built artifact. The fold is a pure function of the repository state, so a reviewer
re-runs the documented command and obtains a byte-for-byte identical manifest id, and a
deployed bytecode that does not match the built artifact fails the package closed with a
fixed reason code. The same seven reason codes, in the same fixed order, are recomputed
on-chain by a reference contract, so a whitelisted firm never has to trust the tool's word
for the verdict. This paper describes the design and the milestone-1 implementation: a
dependency-free Python core and a byte-for-byte Solidity reference, verified by 24 Python
tests and 13 Solidity tests. All work is Arbitrum Sepolia testnet only; no mainnet, no
deployment, no users, and no audit (Section~\ref{sec:limits}).
\end{abstract}

\section{Problem}
An audit is only as deep as the time that survives the setup. When a whitelisted firm
opens a repository cold, it reconstructs a build environment (which compiler version,
which optimizer settings, which library commits), works out how to run the test suite,
and reverse-engineers which on-chain address corresponds to which contract, all before
review begins. Three failure shapes recur. First, the build is \emph{irreproducible}: an
unpinned compiler range or an environment-dependent step means the artifact the firm
builds is not the artifact that was deployed. Second, the evidence is \emph{asserted, not
carried}: the team states that tests pass and that a deployed address holds a given
contract, but nothing in the handoff lets the firm confirm either without redoing the
work. Third, the review is \emph{not re-runnable}: after fixes, a re-audit starts the
same archaeology over again, because the first pass captured nothing a machine can replay.

The cost lands on the subsidy. The Arbitrum Audit Program~\cite{arbaudit} allocates
\$10M in ARB over twelve months to compensate a fixed whitelist of firms for third-party
smart-contract audits on Arbitrum~\cite{arbitrum}. Every hour a firm spends reconstructing
a build is subsidy that bought no findings. The gap is narrow and concrete: there is no
standard, reproducible, evidence-carrying package a project produces \emph{before} the firm
starts, so the firm opens a build that already works, a test suite that already ran with
its counts captured, and a verdict anyone can recompute.

\section{Why Arbitrum}
AuditLane is Arbitrum-specific by construction, and not because of a superficial
integration. Its manifest is content-addressed, and the reproducibility verdict is a pure
function whose seven reason codes and fixed evaluation order are mirrored byte-for-byte in
a reference contract, \texttt{BuildRegistry.sol}. A project commits a manifest id on
Arbitrum Sepolia; anyone recomputes the same id from the published package and checks the
two agree, and the commitment is \emph{write-once} so a claim is immutable once made. The
result is that ``this build reproduced, here is the verdict'' becomes a public, timestamped,
on-chain claim on Arbitrum~\cite{arbdocs}, not a line in a report the reader is asked to
believe. The timing is the program itself: the Audit Program is live, funded for a fixed
window, and routing rolling engagements now, so a tool that reduces setup waste is worth
most while the pipeline is running.

\section{Design}
AuditLane is one pure core, one on-chain reference, and one agent wrapper.

\subsection{The manifest: a pure fold}
The load-bearing component (\texttt{auditlane/manifest.py}) takes a repository spec ---
the pinned toolchain, the source files committed by content hash, the captured test-suite
result, and the on-chain deployments --- and folds it into a canonical payload over which
a single content-addressed \texttt{manifest\_id} is taken. The fold performs no I/O, uses
no clock, and touches no network: every value that varies is passed in, so the digest is a
function of evidence alone~\cite{reproducible}. Sources and contracts are sorted by a
stable key before hashing, so their order in the repository cannot change the id, and the
canonical field order is fixed and mirrored in Solidity. Leaf content is committed as
SHA-256~\cite{sha2} hashes, so the manifest never carries source bytes, only commitments
to them. The consequence is the reproducibility property milestone 1 promises: two runs
over the same repository state produce a byte-identical manifest and the same id, and a
reviewer who clones the repository and re-runs the documented command reproduces the id
and the test counts we publish.

\subsection{The reproducibility check: fixed order, fail closed}
A deterministic \texttt{check} returns the reason codes that apply, in a fixed order; an
empty list is \texttt{REPRODUCIBLE}. The codes, in order, are: an unpinned toolchain
(1), a build that captured no sources (2), a declared test suite that did not pass (3), a
contract whose source is missing from the package (4), a deployed bytecode that does not
match the built artifact (5), and a contract with no resolved on-chain deployment (6).
The check is fail-closed: any unmet condition adds a code, and a green manifest carries
none. Critically, the scanner does \emph{not} synthesise a passing test suite: the counts
come from a captured \texttt{auditlane.tests.json} written after the suite ran, so the
manifest reports the counts that actually ran.

\subsection{The on-chain reference: \texttt{BuildRegistry.sol}}
\texttt{BuildRegistry} recomputes the same verdict on-chain from a package's checkable
shape, with the identical reason-code constants and the identical evaluation order, so the
first-failing code returned on-chain matches the off-chain verdict. \texttt{manifestDigest}
recomputes the id from the canonical payload with Keccak-256~\cite{keccak}, and
\texttt{commit} writes a manifest id once, reverting on a re-commit, so a package's claim
is a public, timestamped, immutable commitment. The two implementations share a common
ABI: the same reason codes, the same order, the same canonical joining.

\subsection{The Arbitrum chain seam}
The on-chain claim AuditLane makes is: ``the contract deployed at this Arbitrum Sepolia
address is the contract in this repository.'' To check that without trusting the tool, the
manifest carries the hash of the deployed runtime bytecode next to the hash of the locally
built artifact, and the check fails closed if they diverge. The seam
(\texttt{auditlane/chain.py}) is one interface with two backends. Offline (the default,
keyless) it returns the deployed bytecode hash from a deterministic in-process fixture,
seeded from the sample's own artifact bytes so the honest case reproduces and a tampered
artifact fails closed. Live (gated on \texttt{AUDITLANE\_RPC\_URL}) it calls
\texttt{eth\_getCode} on Arbitrum Sepolia and refuses any chain id other than 421614. It
is read-only either way and never needs a key.

\subsection{The agent wrapper}
The build runs as a recorded action wrapped in agent-core guardrails
(\texttt{auditlane/agent.py}): the manifest write passes an \texttt{ActionLimiter}, so a
runaway build loop is rate-limited like any other outbound action, and every step (scan,
resolve, reproducibility check, manifest) is journalled in a \texttt{StateStore} so a UI
can replay the activity log. The agent never emits a green manifest a failing check
produced; that is the trust primitive.

\section{Implementation and evaluation}
The core is dependency-free Python: hashing, the canonical payload, and the check are a
few hundred lines with no third-party runtime dependency, so ``offline mode'' can never
degrade into an import error. The reference contract targets solc 0.8.24 and is
self-contained. We evaluate by test, using Foundry~\cite{foundry} for the contract:

\begin{itemize}[leftmargin=1.4em]
\item \textbf{13 Solidity tests} (\texttt{forge test}, solc 0.8.24) on
\texttt{BuildRegistry.sol}: every reason code, the fixed check order, the write-once
commitment, and the Keccak digest.
\item \textbf{24 Python tests} (\texttt{pytest}), keyless and offline: the reproducibility
check for every code and its order, manifest determinism and content-addressing, the
scanner reading a real fixture repository, and the agent recording a build with
fail-closed guardrails.
\end{itemize}

The two suites are built to agree: the \texttt{test\_check\_*} cases are the exact
analogues of the Solidity \texttt{test*Fails} cases, so a green \texttt{pytest} proves the
off-chain core and the on-chain reference produce the same reason code, in the same order,
for the same package. The sharpest case is the tampered deployment: a resolver returns the
wrong deployed bytecode for the sample's address, and the build yields
\texttt{NOT\_REPRODUCIBLE} with \texttt{BYTECODE\_MISMATCH} (code 5), off-chain and
on-chain alike.

\section{Limitations}\label{sec:limits}
This is milestone 1, and the honest scope is narrow.
\begin{itemize}[leftmargin=1.4em]
\item \textbf{No mainnet}, and none planned under this grant. All Web3 targets are Arbitrum
Sepolia testnet only; \texttt{BuildRegistry.sol} is not deployed anywhere yet (testnet
deployment is milestone 2).
\item \textbf{No live-chain test.} The live \texttt{eth\_getCode} path is wired and gated
on \texttt{AUDITLANE\_RPC\_URL}, but the test suite runs the offline fixture, not a real
Arbitrum endpoint.
\item \textbf{No users, no revenue, no audit, no partnership with the Arbitrum Foundation.}
Milestone 3 is a request to be audited through the Audit Program, not a claim that we have
been. The substitute for traction is tested code and a candid scope map, not usage numbers.
\end{itemize}

\section{Conclusion}
AuditLane makes the reproducibility of an Arbitrum project's build a checkable artifact
rather than an assertion: a content-addressed manifest a whitelisted firm opens without
reconstruction, a fail-closed verdict with a fixed reason-code order, and a byte-for-byte
on-chain reference that turns the verdict into a public, recomputable commitment on
Arbitrum. Milestone 1 delivers the cleanroom core and the reference contract, tested. The
durable contribution is a standard, reproducible, evidence-carrying handoff format that
makes every subsidised audit in the program start from a better place.

\bibliographystyle{plain}
\bibliography{references}

\end{document}
